% Part: conditional-logics
% Chapter: minimal-change-semantics
% Section: sphere-models

\documentclass[../../../include/open-logic-section]{subfiles}

\begin{document}

\olfileid{con}{min}{sph}

\olsection{球形模型}

为反事实条件句提供形式语义的一种方法是，将 Lewis 的非形式描述转化为一个数学结构。这样，世界~$w$ 周围的球层就成了世界的集合。由于球层是嵌套的，$w$ 周围的这些世界集合必须按子集关系线性排序。

\begin{defn}
  一个\emph{球形模型}是一个三元组~$\mModel{M} = \tuple{W, O, V}$，
  其中 $W$ 是一个非空的世界集合，$V\colon \PVar \to \Pow{W}$ 是一个赋值，而 $O\colon W \to \Pow{\Pow{W}}$ 为每个世界~$w$ 指派一个\emph{球层系统}~$O_w$。对每个 $w$，$O_w$ 是一个世界集合的集合，并且必须满足：
  \begin{enumerate}
  \item $O_w$ \emph{以 $w$ 为中心}：$\{w\} \in O_w$。
  \item $O_w$ 是\emph{嵌套的}：只要 $S_1, S_2 \in O_w$，就有 $S_1 \subseteq S_2$ 或 $S_2 \subseteq S_1$，即 $O_w$ 按~$\subseteq$ 线性排序。
  \item $O_w$ 对非空并封闭。
  \item $O_w$ 对非空交封闭。
  \end{enumerate}
\end{defn}

$O_w$ 背后的直观是：$w$ “周围”的世界根据它们距离 $w$ 的远近被分层。最内层的球层就是 $w$ 本身，即集合~$\{w\}$：$w$ 比任何其他球层中的世界都更接近 $w$。如果 $S \subsetneq S'$，那么 $S' \setminus S$ 中的世界就比 $S$ 中的世界离 $w$ 更“远”：$S' \setminus S$ 是 $S$ 与 $S'$ 之外世界之间的“层”。特别是，我们必须将球层理解为包含其外表面以内的所有世界；它们不仅仅是单独的层。

\begin{figure}
\begin{center}
\begin{tikzpicture}[layerwidth=1.5,scale=.6]\tiny
  \spheresystem{3}
  \propositionintersect{0}{3}{60}{6}
  \path (0,0) node[world] {$w$};
  \spherepos{-30}{2}{node[world] {$w_2$}}
  \spherepos{220}{2}{node[world] {$w_3$}}
  \spherepos{90}{2}{node[world] {$w_1$}}
  \spherepos[shift={(.1,0)}]{10}{3}{node[world] {$w_5$}}
  \spherepos[shift={(.1,0)}]{-10}{3}{node[world] {$w_6$}}
  \spherepos{225}{3}{node[world] {$w_4$}}
  \spherepos{20}{4}{node[world] {$w_7$}}
  \path (5,0) node[left] {$p$};
\end{tikzpicture}
\caption{球形模型示意图}
\ollabel{fig:sphere-model}
\end{center}
\end{figure}

\olref{fig:sphere-model} 中的示意图对应于这样一个球形模型：$W = \{w, w_1, \dots, w_7\}$，$V(p) = \{w_5, w_6,
w_7\}$。最内层的球层是 $S_1 = \{w\}$。距离 $w$ 最近的世界是 $w_1, w_2, w_3$，所以下一个更大的球层是 $S_2 = \{w, w_1,
w_2, w_3\}$。更远的世界是 $w_4$, $w_5$, $w_6$，因此最外层的球层是 $S_3 = \{w, w_1, \dots, w_6\}$。$w$ 周围的球层系统是 $O_w = \{S_1, S_2, S_3\}$。世界~$w_7$ 不在 $w$ 周围的任何球层中。其中 $p$~为真的最近的世界是 $w_5$ 和~$w_6$，因此最小的 $p$-\allowbreak 容许球层是~$S_3$。

为了在球形模型~$\mModel M$ 中定义公式~$!A$ 在世界~$w$ 处的满足关系，$\mSat{M}{!A}[w]$，我们将模态公式的定义扩展，加入针对 $!B \cif !C$ 的条款：

\begin{defn}
  $\mSat{M}{!B \cif !C}[w]$ 当且仅当要么
  \begin{enumerate}
  \item\ollabel{sphere-vac} 对所有 $u \in \bigcup O_w$，$\mSat/{M}{!B}[u]$，要么
  \item\ollabel{sphere-nonvac} 存在 $S \in O_w$，使得
    \begin{enumerate}
      \item $\mSat{M}{!B}[u]$ 对某个 $u \in S$ 成立，且
      \item 对所有 $v \in S$，要么 $\mSat/{M}{!B}[v]$，要么 $\mSat{M}{!C}[v]$。
    \end{enumerate}
  \end{enumerate}
\end{defn}

根据这个定义，$\mSat{M}{!B \cif !C}[w]$ 当且仅当前件~$!B$ 在 $w$ 周围的所有球层中都为假，或者存在一个球层~$S$，其中 $!B$ 为真，并且实质条件句 $!B \lif !C$ 在该“$!B$-\allowbreak 容许”球层中的所有世界上都为真。注意，定义中并未要求 $S$ 是\emph{最内层的} $!B$-\allowbreak 容许球层，这与直观解释可能给出的预期相反。但是，如果~\olref{sphere-nonvac} 中的条件对某个球层~$S$ 成立，那么它对所有 $S$ 所包含的球层也成立，因此特别地，对最内层的球层也成立。

另请注意，球形模型的定义并不要求一定\emph{存在}一个最内层的 $!B$-\allowbreak 容许球层：我们可能有一个无穷序列 $S_1 \supsetneq S_2 \supsetneq \dots \supsetneq
\{w\}$ 的 $!B$-\allowbreak 容许球层，因而没有最内层的 $!B$-\allowbreak 容许球层。在这种情况下，$\mSat{M}{!B \cif !C}[w]$ 当且仅当对某个~$i$，$!B \lif !C$ 在球层 $S_i$, $S_{i+1}$, \dots 中处处成立。

\end{document}